\begin{theorem}[A sharper five-contact height estimate]\label{thm:quartic-height19}
Keep $S_*$, $\mathcal D$ and $M_5$ from Theorem~\ref{hyp:o8-height}, and put
\[
 B_*=\prod_{p\le76176}p,\qquad
 A_{19}=2^{109}210^{1800},\qquad H_{19}=2A_{19}B_*^{19}+128.
\]
For every $a\in\mathbb Q\setminus\mathcal D$,
\[
 h(a)\le H_{19}M_5(a)^{19},
 \qquad \log_2H_{19}<2093966.
\]
Here $h$ is the absolute logarithmic Weil height.
\end{theorem}
\begin{proof}
We use the geometric construction and compatible representations in the
proof of Theorem~\ref{hyp:o8-height}, not merely its final exponent-$24$
inequality. Write $R=B_*M_5(a)$ and $C=210^{300}$.
The geometric comparison there is
\[
 h(a)\le2\max(0,h_F(A))+128.
\]
All $h_F$ in this proof use the original stable Faltings metric.
The CM estimate in that proof is $h_F(A)<854R$.

In the non-CM case the same construction gives an absolutely simple
quaternionic surface $A/L$, a normal field $L'/\mathbb Q$, and the final
primitive weight-two newform $f$, after the finite twist has been
absorbed, with
\[
 [L:\mathbb Q]\le576,\quad d'=[L':\mathbb Q]\le2304R,\quad
 N\le CR^2,\quad \operatorname{rad}(N)\mid R.
\]
Its two-dimensional compatible representation $r_\ell$ satisfies
\[
 V_\ell(A)|_{G_{L'}}\simeq
       (r_\ell\oplus r_\ell)|_{G_{L'}},\qquad
 (\det r_\ell)^2|_{G_{L'}}=\chi_{\rm cyc}^2|_{G_{L'}}.
\]
These are (LR3)--(LR5) and the determinant calculation in the existing
modularity argument. There is no bounded-degree hypothesis on the
coefficient field of $f$.

Let $\epsilon$ be the nebentypus of this final form. In the Tate-module
convention, with arithmetic Frobenius, $\det r_\ell=\epsilon\chi_{\rm cyc}$;
see \cite[\S12.5, (12.5.1)--(12.5.2) and Remark 12.5.1,
pp.~119--120]{diamond-im}. Equivalently, at a good prime its
